\chapter{Mathematical Foundations}

\section{Trajectory-game model}

Let $i \in [N]$ index agents. At time $t$, each agent state is
$x_t^i \in \mathcal{X}^i$, control is $u_t^i \in \mathcal{U}^i$, and decision
variable is $z_t^i = [x_t^i, u_t^i]$. Over horizon $T$, trajectory
$\mathbf{z}^i = \{z_0^i,\dots,z_{T-1}^i\}$.

A constrained trajectory game for agent $i$ is:
\begin{equation}
\begin{aligned}
\min_{\mathbf{z}^i} \quad & J^i(\mathbf{z}^i,\mathbf{z}^{-i};\theta^i) \\
\text{s.t.}\quad & x_{t+1}^i = f^i(x_t^i,u_t^i), \\
& x_0^i = \hat{x}_0^i, \\
& g^i(\mathbf{z})=0,\; h^i(\mathbf{z})\ge0, \\
& g^s(\mathbf{z})=0,\; h^s(\mathbf{z})\ge0.
\end{aligned}
\label{eq:traj_game_foundation}
\end{equation}

\section{Ordered preferences and lexicographic optimization}

Each agent has an ordered objective list
$\{J_1^i, J_2^i, \dots, J_{K^i}^i\}$, where lower index means higher priority.
GOOP solutions enforce strict precedence through nested optimization
\cite{lee:2025}:
\begin{equation}
\begin{aligned}
\min_{\mathbf{z}_{K^i}^i} \;& J_{K^i}^i(\mathbf{z}_{K^i}^i,\mathbf{z}^{-i}) \\
\text{s.t.}\;& \mathbf{z}_{K^i}^i \in \arg\min_{\mathbf{z}_{K^i-1}^i} J_{K^i-1}^i(\cdot), \\
& \vdots \\
& \mathbf{z}_2^i \in \arg\min_{\mathbf{z}_1^i} J_1^i(\mathbf{z}_1^i,\mathbf{z}^{-i}),
\end{aligned}
\label{eq:goop_nested_foundation}
\end{equation}
plus private and shared feasibility constraints.

\section{Receding-horizon execution}

Long-horizon solves are split into repeated short-horizon solves of length $T$
with execution chunk $T_l \le T$ \cite{christophersen:2007,hall:2022}. At each
step:

\begin{enumerate}
    \item measure current joint state,
    \item solve GOOP over window $T$,
    \item execute first $T_l$ controls,
    \item shift and repeat.
\end{enumerate}

\section{Lexicographic IBR over time}

The baseline method used in this thesis is lexicographic IBR over time
\cite{delasheras:2025}:

\begin{itemize}
    \item decompose multi-agent GOOP into per-agent lexicographic best
    responses,
    \item iterate best responses for a bounded number of rounds $L$,
    \item warm start each step with shifted trajectories from previous solves.
\end{itemize}

This introduces a runtime-quality trade-off controlled by $L$ and makes larger
receding-horizon problems tractable in practice.

\section{Bridge to inverse estimation}

The same model supports inverse questions: if objective primitives are known but
ordering is unknown, estimate ordering by matching predicted and observed
behavior over an identification horizon. This bridge is central to the
\texttt{InverseGoopSolver} extension.

\begin{figure}[tbp]
    \centering
    \fbox{\parbox{0.92\linewidth}{
    \textbf{Placeholder figure: Forward/inverse shared model.}
    Diagram showing that forward solving and inverse estimation share dynamics,
    constraints, and objective primitives, differing mainly in which variables
    are optimized versus inferred.
    }}
    \caption{Shared mathematical core across thesis work packages.}
    \label{fig:shared_core}
\end{figure}
